\RequirePackage{soul}
\RequirePackage{soulpos}
\RequirePackage{tikz}

\makeatletter
% Inline highlights are drawn as a soulpos underlay rather than as a \tcbox. A \tcbox is
% a single unbreakable hbox, so a phrase longer than the space left on the line was
% pushed out whole: it stranded the following words at the right margin, or ran past it
% (a full sentence inside a theorem block overflowed by ~100pt). soulpos calls the
% underlay once per line fragment, so the highlight follows line breaks instead.
\newcommand\thmi@bg{inline-term-bg}
\newcommand\thmi@fg{inline-term-fg}
\newlength\thmi@pad
\setlength\thmi@pad{1.7pt}

\ulposdef{\thmi@underlay}{%
    \begin{tikzpicture}[overlay]
        \draw[
            fill=\thmi@bg,
            draw=\thmi@fg!30!white,
            line width=0.25pt,
            rounded corners=1.6pt
        ]
            (-\thmi@pad, -0.32em) rectangle (\dimexpr\ulwidth+\thmi@pad\relax, 1.02em);
    \end{tikzpicture}%
}

% Breakable inline highlight.
% #1 font declaration (e.g. \bfseries), #2 background color, #3 text color, #4 content.
% Note: the optional argument is a font declaration, not tcolorbox keys.
\NewDocumentCommand{\themeInlineBox}{O{} m m +m}{%
    \begingroup
    \renewcommand\thmi@bg{#2}%
    \renewcommand\thmi@fg{#3}%
    \kern\thmi@pad
    \begingroup\color{#3}#1\thmi@underlay{#4}\endgroup
    \kern\thmi@pad
    \endgroup
}

% Unbreakable variant, kept for content that must not be re-typeset word by word.
% soul cannot scan \detokenize output, and a code identifier has no spaces to break at
% anyway, so \code gains nothing from the breakable path.
\NewDocumentCommand{\themeInlineChip}{O{} m m +m}{%
    \tcbox[
        on line,
        nobeforeafter,
        box align=base,
        boxsep=0pt,
        left=1.7pt,
        right=1.7pt,
        top=0.35pt,
        bottom=0.35pt,
        arc=1.6pt,
        outer arc=1.6pt,
        boxrule=0.25pt,
        colback=#2,
        colframe=#3!30!white,
        colupper=#3,
        #1
    ]{\strut #4}%
}
\makeatother

\newcommand{\key}[1]{\themeInlineBox[\bfseries]{inline-term-bg}{inline-key-fg}{#1}}
\newcommand{\term}[1]{\themeInlineBox{inline-term-bg}{inline-term-fg}{#1}}
\newcommand{\warn}[1]{\themeInlineBox[\bfseries]{inline-todo-bg}{inline-warn-fg}{#1}}
\newcommand{\todo}[1]{\themeInlineBox[\sffamily\footnotesize\bfseries]{inline-todo-bg}{inline-todo-fg}{TODO #1}}
\newcommand{\code}[1]{%
    \themeInlineChip[fontupper=\ttfamily\footnotesize]{inline-code-bg}{inline-code-fg}{\detokenize{#1}}%
}
\makeatletter
\newcommand{\sidenote}[1]{%
    \marginnote{%
        \footnotesize
        \setlength{\parindent}{0pt}%
        \color{sidenote-fg}%
        \textcolor{sidenote-accent}{\rule{1.3pt}{1.15em}}%
        \hspace{3pt}%
        \parbox[t]{0.82\marginparwidth}{\raggedright #1}%
    }%
}
\newcommand{\mnote}[1]{\sidenote{#1}}
\NewDocumentCommand{\chapteroverview}{O{Chapter Overview} +m}{%
    \begin{tcolorbox}[chapteroverview/base,title={#1}]
        #2
    \end{tcolorbox}
}
\makeatother

\newcommand{\lcm}{\operatorname{lcm}}
